\section{Introduction}
\label{sec:introduction}

In a UTXO ledger, a payment spends an output that an earlier transaction created. If the recipient must wait for that transaction to execute publicly, every hop of a payment chain adds a confirmation delay. Suppose Alice pays Bob 100 CAL and Bob forwards the output to Carol before Alice's transaction executes. Bob needs verifiable evidence that the value is backed. If Bob's payment executes first, the ledger consumes an output that does not yet exist, and the missing 100 CAL needs an identified source. The accounting must remain correct whether Alice's transaction arrives before Bob's, after it, or never.

The most direct alternative is to wait for public execution at every hop, which Section~\ref{sec:evaluation} uses as a control. Certificate-based systems already separate recipient-side usability from public processing: FastPay and Zef support low-latency certified payments, and Sui Lutris combines broadcast and consensus paths for different kinds of state access~\cite{fastpay,zef,lutris}. Payment channels reach rapid transfers through channel relationships and their liquidity~\cite{blitz,donner}. Next asks how a native ledger funds a certified payment whose source has not yet executed, with registered organizations, not prefunded channels, standing behind unresolved sources.

Four difficulties are coupled. Certificates circulate before public registration, so accounting that starts from publicly recorded gaps misses existing commitments. A late source must not pay both the guarantor that advanced the value and the recipient who already spent it. A client may hold a source and never submit it. And a certificate is not a promise of execution: $\mathsf{IntentID}=\operatorname{Digest}(\texttt{INTENT},\mathit{Network},\mathit{Subject},\mathit{Nonce})$ is global for a payer on a network, so a malicious user can obtain quorum certificates (QCs) from two organizations for two payments with independent inputs, although public execution accepts only one TxID for that intent. A certificate thus bounds its issuer's liability without guaranteeing that the source executes or repays.

Our insight is that the party answerable for a missing source is already named by the certificate that travels with the output. Next charges the gap to that direct issuer rather than tracing ancestry, bounds every signed certificate's liability by reserve-backed signing capacity, and closes the obligation by a public decision that needs no rewrite of stored history. We organize the design into three contributions.

\noindent\textbf{C1: A layered architecture for one-round continuation.} A member locks inputs and reserves signing capacity atomically before signing, and one parallel round yields a TXCer. The recipient verifies and records the output atomically and can request a successor at once; the request carries the certificates of its direct inputs, so for fixed shapes its evidence does not grow with depth. Replication within the organization and public submission run in the background, and payment principal stays separate from guarantor reserves.

\noindent\textbf{C2: Direct responsibility with source decoupling.} If a certified input is missing, public execution atomically consumes it and records an obligation against the issuer. The successor's block carries the source's certificate, so original signers can resubmit the source from stored approvals. After the deadline, a public decision debits the issuer's reserve without threshold adaptation, and a late source repays it. An unexecuted source keeps its full CAL amount, including change, reserved in signing capacity, and partial approvals also occupy capacity. For four-member organizations and a four-validator committee with at most one static Byzantine member each, we prove unique consumption and that every certificate's outstanding CAL liability, published or not, is funded. For finite, acyclic, source-closed, conflict-free payment sets in which every payment succeeds once, terminal CAL holdings do not depend on arrival order or compensation timing; FUEL is excluded. We do not claim these proofs for general committee sizes.

\noindent\textbf{C3: Authorized historical representation.} Restricted chameleon-hash commands replace the funding reference of a compensated input with the reserve reference, preserving the owner authorization, successor references, and the complete BlockID. The permanent decision is the economic basis, and a successfully executed Revision and Task authorize the exact representation; matching hashes alone do not. Economic closure neither waits for it nor requires it. Its interface benefit is that an honest replica that has verified and executed prefix $H$ can read, in one transaction at the original coordinates $(h,j,k)$, which reserve funded the input. The costs are additional stored records and a local maintenance step per repaired block.

On one M4 Max host, mechanism-internal controls on the frozen build give the following results. With synchronous member commits and unsynchronized wallet stores, a 100-hop chain reaches its last certified receipt in a median 4.557\,s, against 64.476\,s when each hop waits for public execution, and in all 297 fast-chain successors the first send precedes the earliest committee commit of the parent. A fourteen-service mixed workload with organization-paid FUEL completes 84,192 payments in three runs per treatment under client backpressure, reported in Section~\ref{sec:evaluation}; original signers resolve all 24 withheld sources, and compensation and repayment complete with adaptation paused. In 21,600 local hot-page queries without object caches, representation lowers point-read P50 by 7.44--12.32\% against an optimized decision index, mainly through block decoding, at 0.28--2.19\,MiB of added records per block and one 43--214\,ms local maintenance step that excludes network and consensus.